\section{Einleitung}
\label{sec:einleitung}

Aluminium ist ein wichtiger Rohstoff in vielen Industriezweigen, darunter die Automobilbranche, Bau- und Verpackungsindustrie. Die Preisentwicklung von Aluminium ist von großer Bedeutung für Unternehmen, die diesen Rohstoff verwenden, da sie direkte 
Auswirkungen auf die Produktionskosten und die Wettbewerbsfähigkeit haben kann. Beispielsweise können steigende Aluminiumpreise die Gewinnmargen von Unternehmen verringern, während sinkende Preise die Rentabilität verbessern können. Daher ist es für Unternehmen von entscheidender Bedeutung, die zukünftige Preisentwicklung von Aluminium vorherzusagen, um fundierte Entscheidungen treffen zu können. Eine Preisveränderung kann durch verschiedene Faktoren beeinflusst werden, wie z.B. Angebot und Nachfrage, geopolitische Ereignisse, wirtschaftliche Entwicklungen und technologische Fortschritte. Daher ist die Fähigkeit, die zukünftige Preisentwicklung von Aluminium vorherzusagen, von großem Interesse für Unternehmen und Investoren.

Die Schwierigkeit des Projektes liegt darin, dass die Preisentwicklung von Aluminium von einer Vielzahl von Faktoren beeinflusst wird, die oft komplex und miteinander verknüpft sind. Beispielsweise können geopolitische Ereignisse wie Handelskonflikte oder Naturkatastrophen die Produktion und den Handel von Aluminium beeinträchtigen, was zu plötzlichen Preisänderungen führen kann. Darüber hinaus können wirtschaftliche Entwicklungen wie Inflation oder Wechselkursschwankungen die Nachfrage nach Aluminium beeinflussen, was ebenfalls Auswirkungen auf die Preise haben kann. Somit sind nicht nur datenbasierte Komponenten relevant, sondern regelbasierte Ansätze können ebenfalls von Bedeutung sein, um die Preisentwicklung von Aluminium zu modellieren und vorherzusagen.

Die Motivation des Projektes besteht darin, ein Muster in den historischen Preisdaten zu erkennen und diese Informationen zu nutzen, um zukünftige Preisbewegungen vorherzusagen. Mehrere Ansätze können verwendet werden, um dieses Ziel zu erreichen, darunter Zeitreihenanalyse, maschinelles Lernen und statistische Modelle. Der gewählte Ansatz für dieses Projekt ist die Anwendung von Zeitreihenanalyse und maschinellem Lernen, um die Preisentwicklung von Aluminium zu modellieren und vorherzusagen. Hierzu werden zunächst die Rohdaten gesammelt und aufbereitet, um sie für die Analyse nutzbar zu machen. Anschließend werden verschiedene Modelle getestet und verglichen, um die beste Vorhersagegenauigkeit zu erzielen. Das Ziel des Projekts ist es, ein zuverlässiges Modell zu entwickeln, das Unternehmen dabei unterstützt, fundierte Entscheidungen in Bezug auf den Einkauf und die Verwendung von Aluminium zu treffen.

% \begin{itemize}
%     \item Umfang: ca. 1 Seite
%     \item Motivation des Use Cases Aluminiumpreisvorhersage
%     \item Relevanz für Unternehmen und Industrie
%     \item Ziel des Projekts
%     \item Kurze Beschreibung des gewählten Ansatzes
% \end{itemize}
% modelle: Random forrest: XGBoost, LightGBM
% Transformer, LStm, Arima?